% Fundamentos ampliados. Se incluye desde manual_teorico_pedagogico.tex.

\section{Modelado, estabilidad y plano de fases}

\begin{objectivebox}
Construir la base analítica que debe preceder a cualquier diagnóstico de caos:
formulación y escalamiento de modelos, dinámica unidimensional, clasificación lineal
plana, linealización no lineal, nullclines, ciclos límite y bifurcaciones locales.
Cada tema termina con un cálculo reproducible en Toolbox Chaos y con una delimitación
explícita de la evidencia que aporta la gráfica y de la que exige cálculos adicionales.
\end{objectivebox}

\subsection{Modelado, variables de estado y adimensionalización}

\paragraph{Variables de estado y escalas.}
Un \term{modelo dinámico} especifica un estado \(\mathbf{x}\), parámetros
\(\boldsymbol\mu\), una ley de evolución y las hipótesis bajo las cuales esa ley es
válida. Elegir el estado incorpora las variables observables y la
información necesaria para continuar el movimiento. Por ejemplo, la posición de una
masa requiere también su velocidad. La forma de primer orden de una
ecuación de orden \(k\) amplía el estado hasta contener las derivadas requeridas.

La \term{adimensionalización} introduce escalas características y reemplaza
magnitudes con unidad por cocientes sin dimensión. Esto reduce el número de
parámetros independientes, permite comparar experimentos de distinta escala y evita
interpretar como esencial un valor que depende de la unidad elegida. Esta
operación identifica las razones físicas que controlan el régimen.

\paragraph{Ecuación dimensional y transformación.}
Considérese el oscilador de Duffing forzado
\[
 m\ddot x+c\dot x+kx+\alpha x^3=F\cos(\Omega t),
 \qquad m>0,\quad k>0.
\]
Se eligen
\[
 \omega_0=\sqrt{\frac{k}{m}},\qquad
 \tau=\omega_0t,\qquad x=L u,\qquad
 L=\sqrt{\frac{k}{|\alpha|}}.
\]
Como \(\dot x=L\omega_0u'\) y
\(\ddot x=L\omega_0^2u''\), al dividir entre \(kL\) resulta
\[
 u''+2\zeta u'+u+s u^3=f\cos(\nu\tau),
\]
donde las primas representan derivación respecto de \(\tau\) y
\[
 \zeta=\frac{c}{2\sqrt{mk}},\qquad
 s=\operatorname{sgn}(\alpha),\qquad
 f=\frac{F}{kL},\qquad
 \nu=\frac{\Omega}{\omega_0}.
\]
Los cinco parámetros dimensionales \(m,c,k,\alpha,F,\Omega\) quedan organizados en
cuatro combinaciones que separan amortiguamiento, signo de la rigidez cúbica,
intensidad y frecuencia relativa del forzamiento.

\paragraph{Ejemplo de adimensionalización.}
Sea \(m=1\), \(k=4\), \(c=0.8\), \(\alpha=2\), \(F=1\) y \(\Omega=3\), en unidades
compatibles. Entonces
\[
 \omega_0=2,\qquad L=\sqrt2,\qquad
 \zeta=\frac{0.8}{2(2)}=0.2,\qquad
 s=1,
\]
\[
 f=\frac{1}{4\sqrt2}\simeq0.1767767,\qquad
 \nu=\frac32.
\]
La ecuación adimensional queda
\[
 u''+0.4u'+u+u^3=0.1767767\cos(1.5\tau).
\]
Para convertirla en un flujo autónomo se define \(v=u'\) y una fase
\(\phi=\nu\tau\):
\[
 u'=v,\qquad
 v'=-0.4v-u-u^3+0.1767767\cos\phi,\qquad
 \phi'=1.5.
\]
La fase crece sin cota, pero el forzamiento depende de ella módulo \(2\pi\). Una
sección \(\phi=0\pmod{2\pi}\) produce el mapa estroboscópico del oscilador.

\begin{activitybox}{Práctica reproducible: Duffing adimensional}
Abra \gui{Crear sistema}. Use las variables \texttt{u, v, phi}; parámetros
\texttt{zeta=0.2}, \texttt{s=1}, \texttt{f=0.176776695},
\texttt{nu=1.5}; y ecuaciones, en ese orden,
\texttt{v},
\texttt{-2*zeta*v-u-s*u**3+f*cos(phi)} y \texttt{nu}.
Use estado inicial \((0.1,0,0)\), método RK4, paso \(h=0.005\), duración \(200\),
5000 muestras transitorias y umbral de divergencia \(10^6\). Pulse
\gui{Validar definicion} y \gui{Simular y graficar}; después reutilice la trayectoria
en \gui{Retratos 2D} y \gui{Series temporales}. Repita con \(h=0.0025\).
\end{activitybox}

\paragraph{Lectura de las variables escaladas.}
En \(u(\tau)\) se observan transitorio, amplitud y modulación; en el retrato
\((u,v)\), rotación, deformación por el término cúbico y posible acercamiento a un
régimen. La fase representa la variable que hace autónomo el forzamiento y se integra
en el estado ampliado del sistema. Dos trayectorias que se separan tarde pueden
conservar amplitudes, sección estroboscópica y distribución semejantes.

\begin{warningbox}
La validación de hipótesis físicas requiere su contraste experimental o teórico, y la
existencia de un atractor requiere evidencia adicional a una trayectoria acotada. La
comparación entre \(u\) y \(x\) incorpora el factor \(L\), mientras que la fase creciente
se interpreta mediante el forzamiento periódico del oscilador.
\end{warningbox}

\lecturaseccion{Robert C. Hilborn, \emph{Chaos and Nonlinear Dynamics: An
Introduction for Scientists and Engineers}, 2.\textsuperscript{a} ed., capítulo 3,
``An Introduction to Flows'', y apéndice G, ``The Duffing Oscillator''; Tamás Tél y
Márton Gruiz, \emph{Chaotic Dynamics: An Introduction Based on Classical Mechanics},
apéndice A, apartados sobre ecuaciones adimensionales, integración de EDO y programas
de ejemplo; Frank C. Hoppensteadt, \emph{Analysis and Simulation of Chaotic Systems},
2.\textsuperscript{a} ed., sección 6.2, apartados sobre Duffing y simulación.}

\subsection{Flujos de una dimensión, potenciales y formas normales}

\paragraph{Campo vectorial en una dimensión.}
Un flujo autónomo unidimensional tiene la forma
\[
 \dot x=f(x;\mu).
\]
Los equilibrios satisfacen \(f(x^*;\mu)=0\). Entre dos ceros consecutivos, la
continuidad de \(f\) mantiene su signo: \(f>0\) orienta la línea de fase hacia la
derecha y \(f<0\) hacia la izquierda. Si el equilibrio es hiperbólico,
\[
 f_x(x^*;\mu)<0\ \Longrightarrow\ \text{estable},\qquad
 f_x(x^*;\mu)>0\ \Longrightarrow\ \text{inestable}.
\]
La unicidad impide que dos soluciones distintas se crucen y descarta órbitas
periódicas no constantes en la recta.

\paragraph{Representación mediante potencial.}
Toda ecuación escalar puede escribirse formalmente como
\[
 \dot x=-\frac{\mathrm dV}{\mathrm dx},
 \qquad
 V(x)=-\int^x f(s)\,\mathrm ds.
\]
A lo largo de una solución,
\[
 \frac{\mathrm dV}{\mathrm dt}=V'(x)\dot x=-[f(x)]^2\leq0.
\]
Los mínimos locales de \(V\) corresponden a equilibrios estables; los máximos, a
inestables. Su extensión a un sistema general de varias dimensiones requiere una
justificación específica.

\paragraph{Ejemplo algebraico resuelto: saddle--node.}
Para
\[
 \dot x=\mu-x^2
\]
y \(\mu=1/4\), los equilibrios son \(x_\pm=\pm1/2\). Como
\(f_x=-2x\),
\[
 f_x(-1/2)=1>0,\qquad f_x(1/2)=-1<0.
\]
Por tanto, \(x_-=-1/2\) es inestable y \(x_+=1/2\) estable. El potencial
\[
 V(x)=\frac{x^3}{3}-\mu x
\]
tiene un máximo en \(-1/2\) y un mínimo en \(1/2\). Para \(\mu=0\) ambos puntos
colisionan en \(x=0\), donde \(f_x=0\) y la prueba lineal es inconclusa; para
\(\mu<0\) ya no hay equilibrio real.

\paragraph{Las tres formas normales escalares.}
\begin{longtable}{>{\raggedright\arraybackslash}p{0.23\linewidth}
                  >{\raggedright\arraybackslash}p{0.27\linewidth}
                  >{\raggedright\arraybackslash}p{0.40\linewidth}}
\toprule
Mecanismo & Forma normal & Ramas y cambio cualitativo \\
\midrule
\endhead
Saddle--node & \(\dot x=\mu-x^2\) &
Ninguna rama para \(\mu<0\); dos ramas \(x=\pm\sqrt\mu\) para \(\mu>0\).\\
Transcrítica & \(\dot x=\mu x-x^2\) &
Las ramas \(x=0\) y \(x=\mu\) intercambian estabilidad en \(\mu=0\).\\
Pitchfork supercrítica & \(\dot x=\mu x-x^3\) &
\(x=0\) pierde estabilidad y nacen \(x=\pm\sqrt\mu\) estables para \(\mu>0\).\\
\bottomrule
\end{longtable}
Estas ecuaciones son representantes locales después de cambios de coordenadas y
escalas; su uso en un sistema real requiere establecer las hipótesis que conducen a esos polinomios.

\begin{activitybox}{Práctica reproducible: líneas de fase}
En \gui{Crear sistema}, defina una variable \texttt{x}, un parámetro
\texttt{mu} y la ecuación \texttt{mu-x**2}. Para
\(\mu=-0.25,0,0.25\), use los estados iniciales
\(-1,-0.6,-0.4,0,0.4,0.6,1\), paso \(h=0.01\) y duración \(20\).
Guarde las series en escalas separadas. Repita para
\texttt{mu*x-x**2} y \texttt{mu*x-x**3}; antes de simular, dibuje sus ramas y
prediga el destino de cada condición.
\end{activitybox}

\paragraph{Lectura de potenciales y retratos lineales.}
En una serie temporal, la aproximación monótona a un nivel concuerda con una flecha
de la línea de fase; el escape rápido indica que el intervalo y el modelo deben
revisarse. Cerca de una bifurcación aparece \term{ralentización crítica}: la tasa
lineal es pequeña y el transitorio puede ser mucho más largo. Una gráfica contra
\(\mu\) debe distinguir ramas estables e inestables; las ramas inestables requieren
continuación o estados iniciales controlados para su representación.

\begin{warningbox}
Un equilibrio requiere un criterio de convergencia y una rama requiere continuación o
barridos de parámetros que la localicen. En un
punto no hiperbólico \(f'(x^*)=0\), el análisis conserva términos no lineales o
incorpora un estudio de signos.
\end{warningbox}

\lecturaseccion{Steven H. Strogatz, \emph{Nonlinear Dynamics and Chaos with
Applications to Physics, Biology, Chemistry, and Engineering},
2.\textsuperscript{a} ed., secciones 2.1--2.4, 2.7--2.8 y 3.1--3.6; Morris W.
Hirsch, Stephen Smale y Robert L. Devaney, \emph{Differential Equations, Dynamical
Systems, and an Introduction to Chaos}, 3.\textsuperscript{a} ed., capítulo 1,
apartados sobre ecuación logística, cosecha y bifurcaciones.}

\subsection{Sistemas lineales planos y plano traza--determinante}

\paragraph{Matriz del sistema y solución espectral.}
Para el sistema homogéneo
\[
 \dot{\mathbf x}=A\mathbf x,\qquad
 A=\begin{pmatrix}a&b\\c&d\end{pmatrix},
\]
la solución es \(\mathbf x(t)=e^{At}\mathbf x_0\). Los autovalores satisfacen
\[
 p(\lambda)=\lambda^2-\tau\lambda+\Delta=0,\qquad
 \tau=a+d,\quad \Delta=ad-bc.
\]
Así,
\[
 \lambda_{\pm}=\frac{\tau\pm\sqrt{\tau^2-4\Delta}}{2}.
\]
El signo de \(\Delta\) detecta una silla cuando es negativo. Si \(\Delta>0\), el
signo de \(\tau\) decide atracción o repulsión y el discriminante
\(\mathcal D=\tau^2-4\Delta\) separa autovalores reales de complejos.

\paragraph{Clasificación algebraica.}
\begin{longtable}{>{\raggedright\arraybackslash}p{0.27\linewidth}
                  >{\raggedright\arraybackslash}p{0.28\linewidth}
                  >{\raggedright\arraybackslash}p{0.35\linewidth}}
\toprule
Condición & Tipo hiperbólico & Lectura local \\
\midrule
\endhead
\(\Delta<0\) & Silla & Una dirección entra y otra sale.\\
\(\Delta>0,\ \mathcal D>0,\ \tau<0\) & Nodo estable & Dos tasas reales de
contracción.\\
\(\Delta>0,\ \mathcal D>0,\ \tau>0\) & Nodo inestable & Dos tasas reales de
expansión.\\
\(\Delta>0,\ \mathcal D<0,\ \tau<0\) & Foco estable & Rotación con contracción.\\
\(\Delta>0,\ \mathcal D<0,\ \tau>0\) & Foco inestable & Rotación con expansión.\\
\(\tau=0,\ \Delta>0\) & Centro lineal & Órbitas periódicas neutrales en el sistema
lineal.\\
\bottomrule
\end{longtable}
Las fronteras \(\Delta=0\) y \(\mathcal D=0\) contienen casos degenerados; la recta
\(\tau=0\) requiere especial cuidado en problemas no lineales.

\paragraph{Ejemplo de clasificación por traza y determinante.}
Sea
\[
 A=\begin{pmatrix}-1&2\\-2&-1\end{pmatrix}.
\]
Entonces \(\tau=-2\), \(\Delta=5\) y
\(\mathcal D=4-20=-16\), de modo que
\(\lambda_\pm=-1\pm2i\): el origen es un foco estable. Además,
\[
 \frac{\mathrm d}{\mathrm dt}(x^2+y^2)
 =2x(-x+2y)+2y(-2x-y)=-2(x^2+y^2),
\]
por lo que \(r(t)=r_0e^{-t}\). La velocidad angular es
\[
 \dot\theta=\frac{x\dot y-y\dot x}{x^2+y^2}=-2.
\]
La órbita gira en sentido horario, con periodo angular \(\pi\), mientras su radio
decrece exponencialmente. Este cálculo explica tanto la forma como la orientación.

\begin{activitybox}{Práctica reproducible: atlas de sistemas lineales}
En \gui{Crear sistema}, use variables \texttt{x, y} y ecuaciones \texttt{a*x+b*y},
\texttt{c*x+d*y}. Ejecute: silla \((a,b,c,d)=(1,0,0,-1)\); nodo estable
\((-1,0,0,-2)\); centro \((0,-1,1,0)\); y foco estable
\((-1,2,-2,-1)\). Use condiciones \((1,0)\), \((0,1)\), \((1,1)\) y
\((-1,1)\), RK4, \(h=0.01\), \(T=15\). Reutilice cada trayectoria en
\gui{Retratos 2D} y verifique sentido de giro y tasas.
\end{activitybox}

\paragraph{Lectura del plano traza--determinante.}
En un nodo, las órbitas se alinean con la dirección propia de contracción más lenta;
en una silla, salvo sobre la variedad estable, terminan dominadas por la dirección
inestable. En un foco, el espacio entre vueltas cambia exponencialmente. Una escala
automática puede ocultar esa contracción: conviene fijar límites comparables o usar
también \(r(t)\).

\begin{warningbox}
El plano traza--determinante clasifica matrices reales \(2\times2\). El análisis de
sistemas de otra dimensión requiere herramientas acordes con su Jacobiano. Un centro
lineal requiere términos de orden superior para determinar si aparece una espiral débil
o un ciclo límite, y la orientación de giro se obtiene del campo completo además de
\(\tau\) y \(\Delta\).
\end{warningbox}

\lecturaseccion{Morris W. Hirsch, Stephen Smale y Robert L. Devaney,
\emph{Differential Equations, Dynamical Systems, and an Introduction to Chaos},
3.\textsuperscript{a} ed., capítulos 2--4, con énfasis en retratos y clasificación
traza--determinante; Steven H. Strogatz, \emph{Nonlinear Dynamics and Chaos with
Applications to Physics, Biology, Chemistry, and Engineering},
2.\textsuperscript{a} ed., secciones 5.1--5.2.}

\subsection{Jacobiano, linealización y nullclines}

\paragraph{Jacobiano, linealización y nullclines.}
Para
\[
 \dot x=f(x,y),\qquad \dot y=g(x,y),
\]
las \term{nullclines} son los conjuntos \(f=0\) y \(g=0\). Sobre la primera, el
vector es vertical; sobre la segunda, horizontal. Sus intersecciones son equilibrios.
El \term{Jacobiano}
\[
 J(x,y)=
 \begin{pmatrix}
 f_x&f_y\\g_x&g_y
 \end{pmatrix}
\]
produce la aproximación
\(\dot{\boldsymbol\xi}=J(x^*,y^*)\boldsymbol\xi+O(\|\boldsymbol\xi\|^2)\).
Si el equilibrio es hiperbólico, la clasificación lineal conserva la estructura
local cualitativa; si algún autovalor tiene parte real cero, se requieren criterios adicionales.

\paragraph{Ejemplo de Jacobiano y puntos críticos.}
Considérese
\[
 \dot x=x(1-x-y),\qquad
 \dot y=y\left(x-\frac12\right).
\]
Las nullclines son
\[
 x=0\ \text{o}\ y=1-x,\qquad
 y=0\ \text{o}\ x=\frac12.
\]
Sus intersecciones relevantes son
\[
 E_0=(0,0),\qquad E_1=(1,0),\qquad E_c=(1/2,1/2).
\]
El Jacobiano es
\[
 J(x,y)=
 \begin{pmatrix}
 1-2x-y&-x\\
 y&x-\tfrac12
 \end{pmatrix}.
\]
En \(E_0\), \(J=\operatorname{diag}(1,-1/2)\): es una silla. En \(E_1\),
\[
 J(E_1)=\begin{pmatrix}-1&-1\\0&1/2\end{pmatrix},
\]
también una silla. En \(E_c\),
\[
 J(E_c)=
 \begin{pmatrix}-1/2&-1/2\\1/2&0\end{pmatrix},
\quad
 \tau=-\frac12,\quad\Delta=\frac14,\quad
 \mathcal D=-\frac34.
\]
Los autovalores son
\[
 \lambda_\pm=-\frac14\pm i\frac{\sqrt3}{4},
\]
así que el equilibrio de coexistencia es un foco estable local.

\paragraph{Campo de signos.}
En el primer cuadrante, \(x\) aumenta debajo de \(y=1-x\) y disminuye encima;
\(y\) aumenta a la derecha de \(x=1/2\) y disminuye a la izquierda. Esta tabla de
signos predice el giro alrededor de \(E_c\) antes de integrar. También muestra que los
ejes son invariantes porque los factores \(x\) y \(y\) impiden cruzarlos desde datos
positivos.

\begin{activitybox}{Práctica reproducible: nullclines y linealización}
En \gui{Crear sistema}, use variables \texttt{x, y} y ecuaciones
\texttt{x*(1-x-y)} y \texttt{y*(x-0.5)}. Simule con RK4,
\(h=0.01\), \(T=60\), desde \((0.1,0.1)\), \((0.9,0.1)\), \((0.2,0.8)\),
\((0.8,0.8)\), \((1.4,0.2)\) y \((0.4,1.2)\). En una copia del retrato, dibuje las
cuatro nullclines y marque los tres equilibrios. Compare la orientación numérica con
la tabla de signos y repita una trayectoria con \(h/2\).
\end{activitybox}

\paragraph{Lectura del campo y las nullclines.}
Las órbitas cruzan una nullcline con tangente vertical u horizontal, salvo en un
equilibrio. Cerca de \(E_c\), la espiral observada debe contraerse con tasa compatible
con \(\Re\lambda=-1/4\); lejos de él, la curvatura depende del campo completo. Las
separatrices de las sillas organizan regiones del plano; una trayectoria
numérica genérica aproxima esa variedad unidimensional con una precisión dependiente del integrador.

\begin{warningbox}
Una nullcline localiza donde una componente de la velocidad se anula; la trayectoria y
las barreras invariantes se determinan a partir del campo completo. El Jacobiano informa
en una vecindad del equilibrio, mientras que la descripción de cuencas globales y coexistencia exige
la evidencia correspondiente sobre ciclos. Cuando el estado simulado abandona el dominio físico,
la revisión del modelo incorpora esa salida y la conserva visible en la gráfica.
\end{warningbox}

\lecturaseccion{Morris W. Hirsch, Stephen Smale y Robert L. Devaney,
\emph{Differential Equations, Dynamical Systems, and an Introduction to Chaos},
3.\textsuperscript{a} ed., capítulo 8, secciones 8.1--8.5, y capítulo 9, apartados
sobre nullclines, estabilidad, sistemas gradiente y hamiltonianos; Steven H.
Strogatz, \emph{Nonlinear Dynamics and Chaos with Applications to Physics, Biology,
Chemistry, and Engineering}, 2.\textsuperscript{a} ed., secciones 6.1--6.3 y 6.8.}

\subsection{Osciladores, ciclos límite y Poincaré--Bendixson}

\paragraph{Ciclos límite y osciladores planos.}
Una \term{órbita periódica} satisface
\(\mathbf x(t+T)=\mathbf x(t)\) para algún \(T>0\). Un \term{ciclo límite} es una
órbita periódica aislada: las órbitas cercanas pueden aproximarse a ella o alejarse.
Un centro lineal, en cambio, posee una familia continua de órbitas cerradas y ninguna
es aislada.

El teorema de Poincaré--Bendixson, en una formulación útil, afirma que si una órbita
de un flujo plano permanece en una región compacta, su conjunto límite no contiene
equilibrios y se satisfacen las hipótesis regulares, entonces dicho conjunto es una
órbita periódica. El teorema delimita los comportamientos posibles al excluir caos determinista sostenido
en flujos autónomos suaves de dos dimensiones, pero exige demostrar acotación y
ausencia de equilibrios en el conjunto límite.

\paragraph{Ejemplo algebraico resuelto: normal radial.}
Considérese
\[
 \dot x=(1-x^2-y^2)x-y,\qquad
 \dot y=x+(1-x^2-y^2)y.
\]
Con \(r^2=x^2+y^2\),
\[
 \dot r=\frac{x\dot x+y\dot y}{r}=r(1-r^2),\qquad
 \dot\theta=\frac{x\dot y-y\dot x}{r^2}=1.
\]
Los radios equilibrados son \(r=0\) y \(r=1\). Como
\[
 g'(r)=1-3r^2,\qquad g'(0)=1,\qquad g'(1)=-2,
\]
el origen es inestable y el círculo \(r=1\) es un ciclo límite estable. El periodo es
\(T=2\pi\). Para una perturbación radial pequeña \(\eta=r-1\),
\(\dot\eta\simeq-2\eta\), de modo que el multiplicador aproximado de retorno es
\[
 M=e^{-2T}=e^{-4\pi}<1.
\]
Esto cuantifica la atracción transversal y la distingue de la contracción a un punto.

\paragraph{Exclusión de ciclos.}
Si la divergencia
\[
 \nabla\cdot\mathbf f=f_x+g_y
\]
mantiene signo estricto en una región simplemente conexa, el criterio de
Bendixson--Dulac excluye órbitas periódicas contenidas en ella. Un signo que cambia
requiere un criterio adicional para decidir la existencia de un ciclo.

\begin{activitybox}{Práctica reproducible: ciclo estable y periodo}
En \gui{Crear sistema}, use variables \texttt{x, y} y ecuaciones
\texttt{(1-x**2-y**2)*x-y} y
\texttt{x+(1-x**2-y**2)*y}. Elija puntos con radios \(0.2,0.8,1.2,2\), por ejemplo
\((0.2,0)\), \((0.8,0)\), \((1.2,0)\), \((2,0)\). Integre con RK4,
\(h=0.01\), \(T=30\). En \gui{Series temporales}, estime el periodo con diez máximos;
en \gui{Retratos 2D}, compruebe acercamiento desde dentro y fuera. Repita con
\(h=0.005\) y compare periodo y radio medio de las últimas cinco vueltas.
\end{activitybox}

\paragraph{Lectura de ciclos límite.}
Una espiral que se acumula sobre una curva cerrada sugiere atracción orbital. La serie
temporal permite comprobar amplitud y periodo; el retrato muestra aislamiento. Para
una evidencia más directa se fija una sección transversal y se grafica
\(r_{n+1}\) contra \(r_n\): un cruce con la diagonal y pendiente de módulo menor que
uno corresponde a un punto fijo estable del mapa de retorno.

\begin{warningbox}
Una curva cerrada trazada con pocas muestras puede ser un polígono numérico. Observar
varias vueltas exige verificar aislamiento y estabilidad. Poincaré--Bendixson se
aplica directamente a mapas, flujos tridimensionales, campos discontinuos o
proyecciones planas de una órbita 3D.
\end{warningbox}

\lecturaseccion{Steven H. Strogatz, \emph{Nonlinear Dynamics and Chaos with
Applications to Physics, Biology, Chemistry, and Engineering},
2.\textsuperscript{a} ed., secciones 7.2, 7.3, 7.4, 7.5, 7.6 y 8.7; Morris W.
Hirsch, Stephen Smale y Robert L. Devaney, \emph{Differential Equations, Dynamical
Systems, and an Introduction to Chaos}, 3.\textsuperscript{a} ed., capítulo 10,
apartados sobre conjuntos límite, mapas de Poincaré y teorema de
Poincaré--Bendixson, y capítulo 12, apartados sobre Liénard y van der Pol.}

\subsection{Bifurcaciones locales y forma normal de Hopf}

\paragraph{Parámetro y cambio de estabilidad.}
Una \term{bifurcación} es un cambio cualitativo de la dinámica al variar un
parámetro. En una bifurcación local de equilibrio, la linealización pierde
hiperbolicidad: un autovalor real cruza cero o un par complejo conjugado cruza el eje
imaginario. La condición espectral señala el posible punto de bifurcación; las condiciones de
transversalidad y no degeneración deciden si aparece la forma normal genérica.

\paragraph{Mecanismos de bifurcación local.}
\begin{longtable}{>{\raggedright\arraybackslash}p{0.25\linewidth}
                  >{\raggedright\arraybackslash}p{0.30\linewidth}
                  >{\raggedright\arraybackslash}p{0.35\linewidth}}
\toprule
Bifurcación & Firma lineal & Consecuencia de la forma normal \\
\midrule
\endhead
Saddle--node & Un autovalor simple pasa por \(0\) &
Creación o destrucción de dos equilibrios.\\
Transcrítica & Un autovalor pasa por \(0\), con ramas preexistentes &
Dos ramas intercambian estabilidad.\\
Pitchfork & Autovalor \(0\) y simetría apropiada &
Una rama simétrica se divide o se reúne.\\
Hopf & \(\lambda_{1,2}(\mu)=\alpha(\mu)\pm i\omega(\mu)\),
\(\alpha(0)=0\), \(\omega(0)\ne0\) &
Nace o muere un ciclo pequeño; su estabilidad depende del término cúbico.\\
\bottomrule
\end{longtable}

\paragraph{Ejemplo algebraico resuelto: Hopf supercrítica.}
La ecuación compleja
\[
 \dot z=(\mu+i\omega)z-|z|^2z,\qquad \omega>0,
\]
equivale, con \(z=re^{i\theta}\), a
\[
 \dot r=\mu r-r^3,\qquad \dot\theta=\omega.
\]
En el origen, los autovalores son \(\mu\pm i\omega\). Para \(\mu<0\), el origen es
estable y no hay radio positivo equilibrado. Para \(\mu>0\), el origen es inestable
y aparece
\[
 r_*=\sqrt{\mu},\qquad T=\frac{2\pi}{\omega}.
\]
Como \(\partial_r(\mu r-r^3)|_{r_*}=\mu-3\mu=-2\mu<0\), el ciclo es estable. Si
\(\mu=0.09\) y \(\omega=2\), entonces \(r_*=0.3\), \(T=\pi\) y el multiplicador
radial lineal por vuelta es
\[
 M=e^{-2\mu T}=e^{-0.18\pi}\simeq0.568.
\]
La amplitud sigue la ley de escala \(\sqrt\mu\) respecto de \(\mu\).

\begin{activitybox}{Práctica reproducible: cruce de Hopf}
En \gui{Crear sistema}, use variables \texttt{x, y}, parámetros
\texttt{mu} y \texttt{omega=2}, y ecuaciones
\texttt{mu*x-omega*y-(x**2+y**2)*x} y
\texttt{omega*x+mu*y-(x**2+y**2)*y}. Ejecute
\(\mu=-0.09,0,0.09,0.25\) desde \((0.05,0)\) y \((0.8,0)\), con RK4,
\(h=0.005\), \(T=80\). Estime en cada caso el radio tardío y el periodo. Grafique
\(r^2\) contra \(\mu\): la predicción exacta de la forma normal es
\(r^2=\mu\) para \(\mu>0\).
\end{activitybox}

\begin{figure}[H]
\centering
\includegraphics[width=0.78\linewidth]{hopf_bifurcation_example.png}
\caption{Ejemplo generado para Toolbox Chaos de un cambio tipo Hopf. La lectura
correcta combina la estabilidad del equilibrio, la amplitud de la órbita y el
parámetro; la identificación de condiciones de genericidad requiere proyecciones y pruebas adicionales.}
\label{fig:hopf-fundamentos-ampliados}
\end{figure}

\paragraph{Interpretación de la gráfica de bifurcación.}
En un diagrama amplitud--parámetro, la rama \(r=0\) cambia de estabilidad en
\(\mu=0\) y la rama \(r=\sqrt\mu\) nace con tangente vertical. En una simulación, el
tiempo de relajación crece cerca de cero porque la tasa \(|2\mu|\) disminuye. Si el
transitorio se mantiene fijo durante un barrido, ese retraso puede deformar la rama o
hacer parecer que existe histéresis.

\begin{warningbox}
Un par de autovalores con parte real próxima a cero motiva la verificación de una
bifurcación de Hopf mediante cruce transversal, no degeneración y presencia
del ciclo. Un barrido en ambas direcciones distingue transitorio, coexistencia e
histéresis; la estabilidad de una rama requiere un criterio adicional a los estados visitados.
\end{warningbox}

\lecturaseccion{Steven H. Strogatz, \emph{Nonlinear Dynamics and Chaos with
Applications to Physics, Biology, Chemistry, and Engineering},
2.\textsuperscript{a} ed., secciones 3.1--3.6 y 8.1--8.4; John Guckenheimer y Philip
Holmes, \emph{Nonlinear Oscillations, Dynamical Systems, and Bifurcations of Vector
Fields}, capítulo 3, apartados sobre variedades centrales, formas normales y
bifurcaciones locales; Morris W. Hirsch, Stephen Smale y Robert L. Devaney,
\emph{Differential Equations, Dynamical Systems, and an Introduction to Chaos},
3.\textsuperscript{a} ed., sección 8.5, ``Bifurcations'', y capítulo 12, apartado
sobre bifurcación de Hopf.}

\subsection*{Criterios de análisis para sistemas planos}
\addcontentsline{toc}{subsection}{Síntesis operativa de los fundamentos}

\begin{evidencebox}{Cadena mínima antes de hablar de caos}
Primero declare el estado y escale el modelo. Después determine equilibrios,
nullclines y Jacobiano; construya una predicción local y un retrato global; compruebe
el integrador y retire el transitorio. Con esa cadena establecida, calcule
bifurcaciones, exponentes, espectros o dimensiones. Esta secuencia permite
distinguir un mecanismo dinámico de un artefacto de modelado o de discretización.
\end{evidencebox}

\begin{activitybox}{Problema integrador de fundamentos}
Elija un sistema plano propio. Entregue: variables y unidades; forma adimensional;
equilibrios; Jacobiano; clasificación traza--determinante; nullclines; regiones de
signo; una afirmación sobre ciclos basada en Poincaré--Bendixson o Bendixson--Dulac;
cuatro trayectorias reproducibles en \gui{Crear sistema}; y una comparación con
\(h/2\). Termine con dos columnas: ``evidencia disponible'' y ``evidencia adicional requerida''.
\end{activitybox}
